\chapter{Evaluation}
\label{chap:evaluation}

\section{Datasets and experimental splits}

The KITTI odometry runs on sequences 00--10 use the stereo camera streams only.
They do not use inertial measurements and must therefore be reported as a
stereo ORB-SLAM3 baseline, not as stereo-inertial results.

\section{Baselines}

The primary offline baseline is ORB-SLAM3 in stereo mode. Any later
stereo-inertial experiment will be reported as a separate sensor configuration
and will not be pooled with the KITTI stereo results.

\section{Accuracy and recovery metrics}

\section{Latency metrics}

Detector latency was measured after one warm-up prediction using batch size
one, an input size of 640 pixels and CUDA device 0. The same locked manifest of
240 image paths was supplied to COCO PyTorch, Mapillary PyTorch and Mapillary
TensorRT models. Median and 95th-percentile preprocessing, inference and
postprocessing latency were recorded. Accuracy was evaluated separately on all
2,000 labelled images in the converted Mapillary validation split using box and
mask precision, recall, mAP50 and mAP50--95. The COCO model was excluded from
the Mapillary accuracy comparison because its class schema is incompatible;
raw detections on unlabeled KITTI, CARLA and Oxford images were treated only as
qualitative domain-transfer evidence.

\section{Ablation protocol}

\section{Software testing}

\section{Threats to validity}

The detector latency results are specific to the tested RTX 4070, CUDA and
TensorRT environment and exclude ROS queueing, OCR, retrieval and ORB-SLAM3
processing. Mapillary validation measures in-domain segmentation accuracy;
the external datasets were not labelled in the 21-class HumanSLAM schema.
Consequently, external detection counts cannot establish precision, recall or
SLAM improvement. The final localisation claim must instead be supported by
matched ORB-SLAM3/HumanSLAM trajectory, recovery and retrieval experiments.
